\documentclass{article}
\usepackage{graphicx} 
\usepackage{cite}
\usepackage{amsmath,amssymb,amsfonts}
\usepackage{algorithmic}
\usepackage{textcomp}
\usepackage{xcolor}
\usepackage{booktabs}
\usepackage{arydshln}
\usepackage[a4paper, margin=1in]{geometry}
\usepackage{algorithm}
\usepackage{algorithmic}
\usepackage{placeins}
\usepackage{float}
\usepackage{subcaption}

\title{Laboratory 4}
\author{Víctor Brull Rico \and Guillem Martínez Sambola}
\date{May 18, 2026}

\begin{document}

\maketitle

\section{Introduction}

% WHAT TO WRITE:
%   - 1–2 paragraphs giving context: why face detection matters,
%     and that two complementary approaches are studied here.
%   - Mention that Viola–Jones is a classical, hand-crafted
%     feature method, while MTCNN is a modern deep-learning one.
%   - End with a one-sentence outline of the report structure.
%
% SUGGESTED TEXT (fill in / adapt):
 
Face detection is a fundamental problem in computer vision with
applications ranging from user authentication to autonomous systems.
In this laboratory we study two representative detection frameworks:
the classical Viola--Jones method~\cite{viola2004robust}, which
achieves real-time performance through a cascade of Haar-like feature
classifiers trained with AdaBoost, and the Multi-task Cascaded
Convolutional Network (MTCNN)~\cite{zhang2016joint}, a deep-learning
pipeline that jointly performs face detection and facial landmark
alignment.
 
The report is organized as follows. Section~\ref{sec:Methods}
describes the theoretical background of both methods.
Section~\ref{sec:Experiments} presents the implementation details and
experimental results. Section~\ref{sec:Conclusion} summarizes our
findings and conclusions.

\section{Methods}
\label{sec:Methods}

\subsection{Viola--Jones Face Detection}
% WHAT TO WRITE:
%   Summarize the three pillars of the method (cite [1] for each).
 
The Viola--Jones detector~\cite{viola2004robust} is built on three
key contributions.
 
\subsubsection{Integral Image}
% WHAT TO WRITE:
%   - Explain what the integral image is, reproduce Eqs.(1)-(2)
%     from the paper (the recurrences for s and ii), and state
%     why it allows any rectangular sum to be computed in O(1).
 
The \emph{integral image} $ii(x,y)$ at location $(x,y)$ is the sum
of all pixel values above and to the left of $(x,y)$:
\begin{equation}
    ii(x,y) = \sum_{x'\le x,\, y'\le y} i(x',y').
\end{equation}
It is computed in a single pass using the pair of recurrences
\begin{align}
    s(x,y)  &= s(x,y-1) + i(x,y), \label{eq:row_sum}\\
    ii(x,y) &= ii(x-1,y) + s(x,y), \label{eq:int_img}
\end{align}
where $s(x,-1)=0$ and $ii(-1,y)=0$~\cite{viola2004robust}.
Once computed, any rectangular pixel sum requires only four array
references, enabling Haar-like features to be evaluated at arbitrary
scale and position in constant time.
 
\subsubsection{Haar-like Features and AdaBoost}
% WHAT TO WRITE:
%   - Briefly describe the three feature types (two-rect,
%     three-rect, four-rect).
%   - Explain that AdaBoost selects a small number of these
%     features as weak classifiers and combines them into a
%     strong classifier. Give the weak-classifier formula h(x,f,p,θ).
 
Three types of rectangle features are used: two-rectangle features
(difference of horizontally or vertically adjacent regions),
three-rectangle features (center region minus two outer regions), and
four-rectangle features (difference between diagonal pairs). With a
$24\times 24$ detection window the exhaustive set contains over
$160{,}000$ features~\cite{viola2004robust}.
 
AdaBoost selects a small subset of these features as \emph{weak
classifiers}. Each weak classifier $h(x,f,p,\theta)$ depends on a
single feature $f$, a polarity $p$, and a threshold $\theta$:
\begin{equation}
    h(x,f,p,\theta)=
    \begin{cases}
        1 & \text{if } pf(x) < p\theta,\\
        0 & \text{otherwise.}
    \end{cases}
\end{equation}
The final strong classifier is a weighted majority vote of the $T$
selected weak classifiers~\cite{viola2004robust}.
 
\subsubsection{Attentional Cascade}
% WHAT TO WRITE:
%   - Explain the cascade structure: stages of increasing
%     complexity, each trained to achieve near-100 % detection
%     with a moderate false-positive rate.
%   - State the overall false-positive and detection-rate products.
%   - Mention that sub-windows rejected at any stage are
%     immediately discarded, which gives speed.
 
The cascade arranges strong classifiers in a sequence. A sub-window
must pass every stage to be labelled as a face candidate; otherwise
it is immediately discarded. Early, lightweight classifiers reject
the majority of background windows, so only a small fraction of
sub-windows reaches the expensive later stages. The final
Viola--Jones cascade uses 38 layers and 6060 features in total,
yet evaluates on average only 8 features per sub-window on a
natural image~\cite{viola2004robust}.
 
\subsection{MTCNN: Multi-task Cascaded Convolutional Networks}
% WHAT TO WRITE:
%   - Summarize the three-stage pipeline (P-Net, R-Net, O-Net)
%     from the paper [2].
%   - Briefly explain the three loss functions (classification,
%     bounding-box regression, landmark localization) and how
%     they are weighted in each network.
%   - Mention online hard sample mining.
 
MTCNN~\cite{zhang2016joint} proposes a cascaded deep-learning
framework for joint face detection and alignment. Given an input
image, an image pyramid is built and fed through three successive
convolutional networks:
 
\begin{itemize}
    \item \textbf{P-Net (Proposal Network):} A fully convolutional
    network that densely scans the image pyramid and produces initial
    face candidate windows together with bounding-box regression
    vectors. Non-maximum suppression (NMS) is applied to merge
    highly overlapping candidates.
 
    \item \textbf{R-Net (Refine Network):} Receives all P-Net
    candidates and rejects a large fraction of false positives via
    a more complex CNN, followed by bounding-box calibration and NMS.
 
    \item \textbf{O-Net (Output Network):} A still deeper CNN that
    refines the remaining candidates and additionally predicts the
    positions of five facial landmarks (left/right eye, nose,
    left/right mouth corner).
\end{itemize}
 
Each network is trained with three tasks simultaneously. The
face/non-face classification loss is a cross-entropy term, while
bounding-box regression and facial landmark localization are both
formulated as Euclidean regression losses~\cite{zhang2016joint}.
An \emph{online hard sample mining} strategy selects the top $70\%$
hardest samples in each mini-batch to strengthen training without
manual curation.

\section{Experiments}
\label{sec:Experiments}

\subsection{Task 1 -- Understanding the Viola--Jones Code}
% WHAT TO WRITE:
%   - One paragraph per main function explaining its role,
%     inputs, and outputs. Functions to cover:
%       ObjectDetection.m
%       ConvertHaarcascadeXMLOpenCV.m
%       HaarCascadeObjectDetection.m
%       OneScaleObjectDetection.m
%       GetIntegralImages.m
%       ShowDetectionResult.m
 
\subsubsection{Overview of the Provided Functions}
 
\textbf{ObjectDetection.m} is the top-level function. It accepts a
grayscale image (or filename), a Haar cascade \texttt{.mat} file, and
an options structure. Internally it resizes the image (unless
\texttt{Options.Resize} is \texttt{false}), calls
\texttt{HaarCascadeObjectDetection} at multiple scales controlled by
\texttt{Options.ScaleUpdate}, and returns an $n\times 4$ array of
detected object bounding boxes $[x, y, w, h]$.
 
\textbf{ConvertHaarcascadeXMLOpenCV.m} reads an OpenCV XML cascade
file and converts it into a MATLAB \texttt{.mat} structure, making the
trained OpenCV classifiers directly usable by the MATLAB pipeline.
 
\textbf{HaarCascadeObjectDetection.m} iterates over image scales,
calls \texttt{GetIntegralImages} to build the integral image (and its
squared version for variance normalization), and then calls
\texttt{OneScaleObjectDetection} at each scale to apply the cascade.
 
\textbf{OneScaleObjectDetection.m} slides the detection window over
the image at a fixed scale, evaluates each stage of the cascade, and
collects accepted sub-windows as object hypotheses.
 
\textbf{GetIntegralImages.m} (detailed in Task~2) computes the
integral image and the squared-pixel integral image from the input
grayscale image.
 
\textbf{ShowDetectionResult.m} draws the detected bounding boxes on
the original image for visualization.
 
\subsubsection{Explanation of Lines 35--38 in OneScaleObjectDetection.m}
% WHAT TO WRITE:
%   - Show (or describe) the four lines.
%   - Explain their role in the context of variance normalization
%     as described in Section 5.4 of [1].
%   - The lines typically compute the mean and variance of the
%     current sub-window using the two integral images, then
%     derive a normalization factor used to scale feature responses.
 
Lines 35--38 of \texttt{OneScaleObjectDetection.m} implement the
\emph{variance normalization} step described in Section 5.4 of
\cite{viola2004robust}. The sub-window mean $\mu$ is obtained from
the integral image via four lookups. The second integral image (of
squared pixel values) yields $\sum x^2$; from these two quantities
the variance is computed as
$\sigma^2 = \tfrac{1}{N}\sum x^2 - \mu^2$,
where $N$ is the number of pixels in the window.  The square root
$\sigma$ is then used as a normalization factor: instead of
normalizing the pixels, the feature thresholds are divided by
$\sigma$ before comparison, keeping the computation inside the
integral-image framework at constant cost per sub-window.
 
% ADD HERE: paste the actual four lines as a verbatim/lstlisting block
% and annotate each line with comments.
\begin{verbatim}
% Loop through a classifier tree
for i_tree=1:length(Trees)
    Tree = Trees(i_tree).value;
    % Executed the classifier
    TreeSum=TreeObjectDetection(zeros(size(x)),Tree,Scale,x,y,IntegralImages.ii,StandardDeviation,InverseArea);
    StageSum = StageSum + TreeSum;
end
\end{verbatim}
 
\subsection{Task 2 -- Implementing the Integral Image}
% WHAT TO WRITE:
%   - Reproduce (or paraphrase) the missing code you implemented
%     in GetIntegralImages.m.
%   - Relate each line to Eqs. (1)-(2) of [1].
%   - Show a small sanity-check result (e.g. sum of a known region).
 
The missing code in \texttt{GetIntegralImages.m} implements
Equations~\eqref{eq:row_sum}--\eqref{eq:int_img}. Concretely:
 
% PASTE YOUR MATLAB CODE BELOW (replace the placeholder):
\begin{verbatim}
% --- MISSING CODE (GetIntegralImages.m) ---
[rows, cols] = size(Picture);
IntegralImage        = zeros(rows+1, cols+1);
IntegralImageSquared = zeros(rows+1, cols+1);
 
for y = 1:rows
    s        = 0;
    s_sq     = 0;
    for x = 1:cols
        s    = s    + Picture(y, x);
        s_sq = s_sq + Picture(y, x)^2;
        IntegralImage(y+1, x+1)        = IntegralImage(y, x+1)   + s;
        IntegralImageSquared(y+1, x+1) = IntegralImageSquared(y, x+1) + s_sq;
    end
end
\end{verbatim}
 
The outer loop traverses rows and the inner loop accumulates the
cumulative row sum $s(x,y)$ (cf.\ Eq.~\eqref{eq:row_sum}). Adding
the value from the row above yields $ii(x,y)$
(cf.\ Eq.~\eqref{eq:int_img}). An identical construction on squared
pixel values produces the second integral image needed for variance
normalization.
 
% ADD A FIGURE: show the integral image of one of your test images.
% \begin{figure}[H]
%   \centering
%   \includegraphics[width=0.5\linewidth]{figures/integral_image.png}
%   \caption{Integral image of \texttt{bruce3.jpg}.}
%   \label{fig:integral_image}
% \end{figure}
 
\subsection{Task 2 -- Merging Overlapping Detections}
% WHAT TO WRITE:
%   - Explain the problem: the sliding-window approach at multiple
%     scales triggers multiple overlapping bounding boxes for the
%     same face.
%   - Describe your merging strategy (e.g. grouping overlapping
%     boxes and averaging their corners, as suggested in [1]
%     Section 5.6, or a non-maximum suppression variant).
%   - Paste your auxiliary function (or its key lines).
%   - Show a before/after figure.
 
Because the detector is evaluated at multiple scales and positions,
each face typically produces several overlapping hypotheses.
Following the strategy proposed in Section~5.6 of
\cite{viola2004robust}, we implemented an auxiliary function
\texttt{MergeDetections.m} that (i) partitions detections into
disjoint subsets whose bounding boxes overlap, and (ii) replaces each
subset with a single box whose corners are the arithmetic mean of the
corners in that subset.

Main loop in the MergeDetections.m function:

\begin{verbatim}
for i = 1:N
    if assigned(i), continue; end
    % Grow cluster of overlapping boxes (transitive closure)
    cluster = i;
    k = 1;
    while k <= numel(cluster)
        idx = cluster(k);
        for j = 1:N
            if ~assigned(j) && j ~= idx && ~ismember(j, cluster)
                if bbox_iou(Objects(idx,:), Objects(j,:)) > overlapThresh
                    cluster(end+1) = j;
                end
            end
        end
        k = k + 1;
    end
    assigned(cluster) = true;
    MergedObjects(end+1, :) = mean(Objects(cluster, :), 1);
end
\end{verbatim}

Auxiliary function \texttt{bbox\_iou} computes the Intersection over union
between two bounding boxes, which is used to determine if they belong to the same cluster.

\begin{verbatim}
function o = bbox_iou(a,b)
% a and b are [x y w h]
ax = a(1); ay = a(2); aw = a(3); ah = a(4);
bx = b(1); by = b(2); bw = b(3); bh = b(4);

ix1 = max(ax, bx);
iy1 = max(ay, by);
ix2 = min(ax+aw, bx+bw);
iy2 = min(ay+ah, by+bh);
iw = ix2 - ix1;
ih = iy2 - iy1;
if iw <= 0 || ih <= 0
    o = 0;
    return;
end
inter = iw * ih;
areaA = aw * ah;
areaB = bw * bh;
o = inter / (areaA + areaB - inter);
end
\end{verbatim}
 
% ADD BEFORE/AFTER FIGURES:
% \begin{figure}[H]
%   \centering
%   \begin{subfigure}[b]{0.45\linewidth}
%     \includegraphics[width=\linewidth]{figures/before_merge.png}
%     \caption{Before merging.}
%   \end{subfigure}
%   \hfill
%   \begin{subfigure}[b]{0.45\linewidth}
%     \includegraphics[width=\linewidth]{figures/after_merge.png}
%     \caption{After merging.}
%   \end{subfigure}
%   \caption{Effect of the merging auxiliary function on a test image.}
%   \label{fig:merge}
% \end{figure}
 
\subsection{Task 2 -- Viola--Jones Face Detection Results}
% WHAT TO WRITE:
%   - Choose at least 3–4 test images (vary: group photo,
%     profile face, occluded face, different lighting).
%   - For each image: show the image with bounding boxes,
%     state how many faces are present and how many were
%     detected (true positives, false positives, missed).
%   - Discuss failure cases explicitly (e.g. profile faces,
%     strong backlighting, occlusion) referencing [1] §5.7.2.
 
We tested the Viola--Jones detector using the
\texttt{haarcascade\_frontalface\_alt.xml} cascade on a set of images
with varying conditions.
 
% EXAMPLE — add one paragraph + figure per image:
%
% \textbf{Image 1 (bruce3.jpg).} The image contains X frontal
% faces. The detector found Y bounding boxes before merging and
% Z after merging, correctly identifying W faces with V false
% positives.
%
% \begin{figure}[H]
%   \centering
%   \includegraphics[width=0.6\linewidth]{figures/bruce3_result.png}
%   \caption{Detection result on \texttt{bruce3.jpg}.}
%   \label{fig:bruce3}
% \end{figure}
 
\subsubsection{Failure Cases}
% Discuss at least one clear failure: tilted faces, occlusion,
% dark faces on light background.  Reference [1] Section 5.7.2.
 
As noted in \cite{viola2004robust}, the cascade was trained on
frontal, upright faces. Consistent with this, our tests revealed
failures in three scenarios: (i) faces rotated more than
approximately $\pm 15^\circ$ in-plane, (ii) strongly backlit faces
where the face region is significantly darker than the background,
and (iii) faces with occluded eyes, since the eye-darkness feature
is one of the first selected by AdaBoost.
 
% ADD FIGURE showing a failure case.
 
\subsection{Task 3 -- MTCNN Results}
% WHAT TO WRITE:
%   - Briefly describe the main.m demo: what it does, what
%     parameters can be set (minSize, threshold per stage, etc.).
%   - Test on the same images used for Viola-Jones for
%     a fair comparison.
%   - Report detected faces and facial landmarks (5 points).
%   - Compare qualitatively: where does MTCNN succeed where
%     VJ fails, and vice versa?
 
We ran the MTCNN demo (\texttt{main.m}) on the same set of images
used in Section~\ref{sec:Experiments}.
MTCNN operates on an image pyramid and refines candidates through
P-Net, R-Net, and O-Net sequentially~\cite{zhang2016joint}.
In addition to bounding boxes it outputs five facial landmarks per
detected face.
 
% EXAMPLE per image:
% \begin{figure}[H]
%   \centering
%   \includegraphics[width=0.6\linewidth]{figures/mtcnn_bruce3.png}
%   \caption{MTCNN result on \texttt{bruce3.jpg}: bounding boxes
%            and five landmarks per face.}
%   \label{fig:mtcnn_bruce3}
% \end{figure}
 
\subsubsection{Comparison: Viola--Jones vs.\ MTCNN}
% WHAT TO WRITE:
%   - Make a small table: image / VJ detections / MTCNN detections
%     / ground-truth faces.
%   - Comment on speed (VJ is real-time on CPU; MTCNN is faster
%     on GPU [2] Table II).
%   - Comment on robustness (MTCNN handles pose variation better).
 
Table~\ref{tab:comparison} summarizes the detection counts for both
methods across the tested images.
 
\begin{table}[H]
\centering
\caption{Detection counts: Viola-Jones vs.\ MTCNN.
TP = true positives, FP = false positives, FN = false negatives.}
\label{tab:comparison}
\begin{tabular}{lccccccc}
\toprule
& & \multicolumn{3}{c}{Viola--Jones} & \multicolumn{3}{c}{MTCNN} \\
\cmidrule(lr){3-5}\cmidrule(lr){6-8}
Image & GT faces & TP & FP & FN & TP & FP & FN \\
\midrule
% FILL IN YOUR NUMBERS:
bruce3.jpg   & -- & -- & -- & -- & -- & -- & -- \\
image2.jpg   & -- & -- & -- & -- & -- & -- & -- \\
image3.jpg   & -- & -- & -- & -- & -- & -- & -- \\
\bottomrule
\end{tabular}
\end{table}
 
Overall, MTCNN proved more robust to moderate pose variation and
partial occlusion, consistent with the benchmark results
reported in \cite{zhang2016joint}, while the Viola--Jones detector
remained competitive on frontal, well-lit faces and offered lower
computational requirements without a GPU.

\section{Conclusion}
\label{sec:Conclusion}

% WHAT TO WRITE:
%   - 1 paragraph summarizing what was implemented and learned.
%   - 1 paragraph comparing the two approaches (speed, accuracy,
%     robustness, ease of use).
%   - Optional: one sentence on possible extensions (video, other
%     cascades, fine-tuning MTCNN).
 
In this laboratory we implemented and evaluated two face detection
frameworks. For Viola--Jones~\cite{viola2004robust}, we completed the
integral image computation following Equations~\eqref{eq:row_sum}
and~\eqref{eq:int_img}, explained the variance normalization step in
\texttt{OneScaleObjectDetection.m}, and implemented an auxiliary
function to merge overlapping detections. For
MTCNN~\cite{zhang2016joint}, we ran the provided demo and analysed
the three-stage cascade that jointly detects faces and localizes five
facial landmarks.
 
The classical Viola--Jones method provides fast, real-time detection
on standard CPUs for frontal faces, but degrades significantly with
pose variation and occlusion. MTCNN leverages deep convolutional
features learned from large datasets and jointly optimizes detection
and alignment, yielding markedly better robustness at the cost of
requiring a GPU for real-time performance. Both methods share the
cascade design philosophy: coarse-to-fine processing that discards
easy negatives early to focus computation on promising regions.

\bibliographystyle{IEEEbib}
\bibliography{references}

\end{document}